\section{系统实现}

\subsection{工作流组织}

系统由四类角色组成：规划智能体负责生成分析蓝图，规格化智能体负责把研究意图转换为可执行技术规格，执行智能体负责代码生成与运行，评审智能体负责结果校验与报告组织。四类角色以状态机方式串联，使每轮规划都能留下显式中间状态，便于复验和错误定位。对期刊稿而言，这一工作流不仅服务于单次分析，还服务于跨模式、跨数据集的可比评测。

\subsection{结构化计划表示}

为保证实验模式可比，本文统一使用结构化 JSON 作为计划表示层。每个任务至少包含：任务类型、变量或列引用、方法名称、参数、预期结论、依赖关系和输出文件位置。采用结构化计划而不是长文本说明，有两个直接收益：一是自动化指标可以稳定提取，二是不同模式之间可以在同一评价协议下比较。规划阶段的所有模式，包括单轮基线、双轮迭代模式和外部机制基线，最终都被归一到同一种中间表示，从而避免“不同系统、不同输出格式”导致的比较偏差。

\subsection{工程接口与可追踪性}

当前系统已固定如下工程接口：

\begin{itemize}
  \item 输入层：研究问题、数据预览摘要、图预览摘要、知识开关、预算约束；
  \item 规划层：任务分解、变量引用、方法选择、参数设置；
  \item 执行层：代码脚本、日志、中间结果、结果摘要；
  \item 评估层：结构指标、语义评分、重复实验汇总。
\end{itemize}

这一设计使本文不仅能够比较最终分数，也能比较不同模式如何组织方法、调用证据和扩展任务深度。对于投稿稿件而言，这一点比“系统能否跑通单个案例”更重要，因为它直接决定了方法能否被第三方复验。

\subsection{实验注册表与证据管理}

为避免学位论文与期刊稿相互覆盖，本文将实验配置与证据目录显式拆分。数据集由注册表统一维护，当前包括主数据集、相邻技术子领域稳定性数据集，以及软件工程 Bugzilla 跨领域数据集；实验模式同样由注册表统一维护，覆盖四方案主实验、两组内部消融，以及两类 literature-inspired external mechanism baselines。每个证据池都独立保存原始运行结果、重复实验汇总、统计检验和图表产物，正文只引用相应证据池导出的结果，而不直接混写不同目录中的原始结果。具体目录结构与导出脚本放入补充材料，不在正文中逐项罗列。

\subsection{投稿级复现实验支撑}

期刊稿工作区已经在 \texttt{\detokenize{scripts/jos/}} 中补齐统计检验、表格导出、图形导出、子领域稳定性批量续跑、Bugzilla 数据集物化、盲评材料导出和盲评汇总等脚本。这些脚本的作用不是替代正文论证，而是确保正文、图表与补充材料中的所有数值都能追溯到稳定的证据目录和统一的导出流程。

\subsection{当前实现边界}

本文工作区已经将期刊稿与学位论文分离，并将所有新增产物限制在 \texttt{\detokenize{docs/paper_jos/}}、\texttt{\detokenize{scripts/jos/}} 与 \texttt{\detokenize{outputs/jos_artifacts/}} 之内。当前已完成的核心证据包括内部 90 组重复实验与 30 组外部机制基线；待补的是 5 问题口径下的扩展重复实验、子领域稳定性全量重跑、Bugzilla 跨领域小验证和真实人工盲评回收。这些缺口已经不再是工程接口缺口，而是待回填的增强版证据缺口。
